\chapter{Conclusion}
\label{sec:conclusion}

This thesis set out to answer a question that is both practical and principled: does a sophisticated tabular transformer architecture outperform a well-tuned gradient-boosted decision tree on financial fraud detection, and if not, why? The question matters because the answer shapes model selection decisions in a high-stakes domain and because understanding the \emph{reasons} behind a performance gap is scientifically more valuable than merely reporting it. The experimental programme — spanning two benchmark datasets, six model families, seven imbalance handling strategies, four threshold selection strategies, cross-domain generalisation, SHAP interpretability, and transformer attention diagnostics, all under a leakage-free protocol — provides a structured and replicable body of evidence on which to ground the answer.

% ──────────────────────────────────────────────────────────────────────────
\section{Summary of Contributions}
\label{sec:contributions_summary}
% ──────────────────────────────────────────────────────────────────────────

The six contributions of this thesis build a cumulative case rather than six independent findings.

The foundation is a \textbf{leakage-free unified comparison} (Contribution~1) that establishes, for the first time in the fraud detection literature, a side-by-side evaluation of classical supervised models and the FT-Transformer under an identical protocol — identical datasets, identical metrics, identical hyperparameter tuning budget, and an outer holdout that was never touched during any selection or calibration step. Without this foundation, the subsequent findings would be confounded by the methodological inconsistencies that pervade the published literature.

Built on that foundation, the \textbf{threshold selection study} (Contribution~2) reveals that the choice of decision threshold is not a minor operational detail but a design decision with first-order impact on deployed performance. On BAF Base, where model output probabilities are concentrated far below 0.5, fixing the threshold at the conventional default reduces F$_2$ to near zero for every supervised model; the max-F$_2$ strategy recovers values in the 0.29--0.32 range — an improvement of approximately nine-fold. No fraud detection system should be deployed without an explicit threshold calibration step.

The \textbf{full factorial analysis} (Contribution~3) across seven imbalance handling strategies, five supervised models, and two datasets reveals that strategy effectiveness is neither universal nor model-agnostic. The most important finding is an architecture-specific interaction: SMOTE degrades FT-Transformer performance by 41\% on BAF (PR-AUC 0.180\,$\to$\,0.106) compared to only 14\% for CatBoost (0.180\,$\to$\,0.154). This differential sensitivity exposes a structural weakness of neural architectures under synthetic oversampling that does not apply equally to tree-based models. Class Weights — which modifies the loss function without altering the data distribution — is the safest default across model families when no prior evidence about strategy effectiveness on the target data is available, although individual exceptions exist (e.g., LGBM on ULB loses approximately 17\% PR-AUC under Class Weights, in line with its sensitivity to every non-baseline strategy on that dataset).

The \textbf{cross-domain generalisation study} (Contribution~4) demonstrates that in-domain performance ranking and cross-domain robustness are not the same thing. Random Forest, the weakest of the high-capacity supervised classifiers in-domain (PR-AUC 0.159 on BAF Base, against 0.177 for LGBM and 0.180 for CatBoost), improves on four out of five BAF variants under zero-shot transfer, while the strongest in-domain models (LGBM and CatBoost) lose 24--31\% of their PR-AUC on covariate-shifted variants. This inversion has direct operational implications: when distribution shift is anticipated, deploying a model solely on the basis of in-domain performance is a suboptimal strategy.

The \textbf{SHAP cross-model and cross-variant analysis} (Contribution~5) establishes two complementary results. First, LGBM and CatBoost converge on an identical top-10 feature ranking (Jaccard\,=\,1.00), providing strong evidence that these features carry genuine fraud signal rather than model-specific artefacts. Second, that ranking is perfectly stable across all six BAF datasets (Jaccard\,=\,1.00 for every pairwise comparison), meaning that explanations generated on the training distribution remain valid under label shift and covariate shift. The FT-Transformer learns a partially different feature hierarchy (Jaccard\,=\,0.54 vs.\ tree-based models), emphasising continuous features — a consequence of its tokenisation design rather than superior discriminative ability.

Finally, the \textbf{FT-Transformer attention diagnostics} (Contribution~6) close the loop on the central question. The Foggy Vision pattern (normalised entropy\,=\,0.976) indicates that attention is near-uniformly distributed over all 31 feature tokens, with no dominant pathological pattern. This diagnosis has a specific implication: the FT-Transformer's competitive parity with CatBoost on BAF (both at PR-AUC\,=\,0.180) is genuine — the model has genuinely processed all available information. The shared performance ceiling reflects a limit of the feature space, not a failure of the architecture.

% ──────────────────────────────────────────────────────────────────────────
\section{Answers to Research Questions}
\label{sec:rq_answers}
% ──────────────────────────────────────────────────────────────────────────

\paragraph{RQ1 — Comparative performance under a leakage-free protocol.}
Under an identical leakage-free evaluation protocol, gradient-boosted models hold a consistent advantage on the smaller ULB dataset: CatBoost leads at PR-AUC\,=\,0.885, followed by LGBM (0.874), RF (0.867), and FT-Transformer (0.856), with Logistic Regression substantially lower (0.742) and OCSVM serving only as a weak anomaly-detection baseline (0.335). On the larger BAF Base ($N$\,=\,1M), the FT-Transformer closes the gap entirely, matching CatBoost exactly at PR-AUC\,=\,0.180. The performance gap is therefore dataset-dependent: on small tabular datasets, GBDTs' inductive bias and calibration advantage translate into a meaningful lead; on large-scale data, the transformer's capacity compensates for its structural mismatch with tabular features. In neither setting does the FT-Transformer outperform the best GBDT, and its substantially higher computational cost provides no corresponding performance benefit over CatBoost or LGBM on these benchmarks (Section~\ref{sec:baseline_results}).

\paragraph{RQ2 — Impact of threshold selection strategy.}
Threshold selection has a profound and dataset-dependent impact on operational metrics. On ULB~2013, the gap between the worst and best threshold strategy is moderate — approximately 13--20\% relative improvement in F$_2$. On BAF~Base, the impact is an order of magnitude larger: the fixed $\tau=0.5$ default reduces F$_2$ to near zero for every supervised model, while the max-F$_2$ strategy recovers F$_2$ in the 0.29--0.32 range — a nine-fold improvement for the FT-Transformer (0.035\,$\to$\,0.317). Model rankings are preserved across threshold strategies on ULB but diverge on BAF, where calibration differences between models become more visible once the threshold is optimised. Threshold selection is a first-class design decision that must be calibrated on a held-out validation set and never fixed at 0.5 without justification (Section~\ref{sec:threshold_sensitivity}).

\paragraph{RQ3 — Factorial interaction of imbalance strategies and model type.}
No single imbalance handling strategy universally improves performance across all models and datasets. On BAF~Base, no strategy improves CatBoost or LGBM beyond their baselines by more than 0.01~PR-AUC, confirming that well-tuned GBDTs handle moderate imbalance natively. The critical architecture-specific finding is SMOTE's asymmetric impact: CatBoost loses only 14\% PR-AUC under SMOTE on BAF, while the FT-Transformer loses 41\% — a difference attributable to the neural model's sensitivity to distributional changes in the training data that disrupt gradient-based optimisation. Class Weights is, on balance, the safest default across model families and is the recommended starting point when no prior evidence exists about strategy effectiveness on the target data, with the caveat that LGBM on ULB is an exception (it degrades under every non-baseline strategy on that dataset, including Class Weights) (Section~\ref{sec:transformer_imbalance}).

\paragraph{RQ4 — Cross-domain generalisation under distribution shift.}
Cross-domain robustness does not follow in-domain performance ranking. Variant~II (label shift only) is the easiest transfer target for all models. Variants~III and~V (covariate shift with partial feature changes) cause the largest degradation: LGBM and CatBoost lose 24--31\% of their Base PR-AUC, and the FT-Transformer degrades similarly (28--34\%), providing no robustness advantage over gradient-boosted models. In contrast, Random Forest improves on four out of five variants despite trailing the gradient-boosted models in-domain (PR-AUC 0.159 vs.\ 0.177--0.180) — an outcome attributable to the variance reduction introduced by random feature subsampling, which prevents overfitting to Base-specific feature co-occurrence patterns. When distribution shift is anticipated, cross-domain robustness must be evaluated explicitly and factored into model selection (Section~\ref{sec:cross_domain_results}).

\paragraph{RQ5 — Feature attribution consistency and transformer attention diagnostics.}
LGBM and CatBoost converge on an identical top-10 feature ranking (Jaccard\,=\,1.00) that is perfectly stable across all six BAF datasets, demonstrating that the most informative fraud predictors are invariant to the distributional shifts of the BAF suite. The FT-Transformer learns a partially different hierarchy (Jaccard\,=\,0.54 vs.\ tree-based models), prioritising continuous features as a consequence of its tokenisation design. Attention diagnostics reveal Foggy Vision (normalised entropy\,=\,0.976): attention is distributed near-uniformly across all feature tokens, with no pathological concentration. This is mechanistically important — the FT-Transformer is not failing to identify informative features; it has genuinely found all available signal, and the shared 0.18 PR-AUC ceiling reflects the information limit of the current feature set rather than a model-specific failure (Sections~\ref{sec:shap_results} and~\ref{sec:attention_diagnostics}).

% ──────────────────────────────────────────────────────────────────────────
\section{Limitations}
\label{sec:limitations}
% ──────────────────────────────────────────────────────────────────────────

Several limitations constrain the scope and generalisability of the findings beyond the protocol-level threats to validity discussed in Section~\ref{sec:discussion}.

\textbf{Single deep learning architecture.} The FT-Transformer is the sole representative of the tabular deep learning family. The conclusions about attention behaviour, imbalance sensitivity, and cross-domain robustness apply directly to this architecture and may not generalise to SAINT \cite{somepalli2021saint}, TabTransformer \cite{huang2020tabtransformer}, TabNet, or hybrid architectures. Whether the Foggy Vision attention pattern is inherent to self-attention on tabular data or specific to FT-Transformer's design is an open question.

\textbf{Dataset scope.} The ULB dataset is privacy-constrained — PCA-transformed features preclude domain interpretation and the two-day collection window limits temporal generalisability. The BAF suite is synthetically generated, which ensures experimental control but means fraud patterns may not fully reflect the complexity of real production fraud streams, where label noise, adversarial adaptation, and operational feedback loops are present.

\textbf{Absence of concept drift modelling.} All experiments use a static train/test split. In operational fraud detection, the data distribution shifts continuously as fraudsters adapt to deployed models. The experimental protocol does not capture this dynamic, and performance under temporal drift is not equivalent to performance on the distribution-shifted BAF variants studied here.

\textbf{Hyperparameter optimisation budget.} The 50-trial Optuna budget was applied equally to all models. The FT-Transformer's search space — number of transformer blocks, token dimension, attention heads, FFN multiplier, dropout rates, learning rate, batch size — is substantially larger than that of classical models, which converge more efficiently under TPE sampling. The budget may therefore under-explore the transformer's potential relative to the gradient-boosted models.

\textbf{Cross-dataset incompatibility.} The ULB and BAF feature spaces are incompatible (PCA-anonymised vs.\ semantic), preventing cross-architecture, cross-dataset comparison. The cross-domain study is therefore restricted to the BAF internal transfer setting, which, while controlled, represents a narrower robustness evaluation than matched feature spaces across independent datasets would provide.

% ──────────────────────────────────────────────────────────────────────────
\section{Future Work}
\label{sec:future_work}
% ──────────────────────────────────────────────────────────────────────────

Five directions emerge directly from the findings and limitations of this thesis.

\textbf{Extending attention diagnostics to other tabular transformer architectures.} The Foggy Vision pattern identified in FT-Transformer raises the question of whether near-uniform attention is an inherent property of self-attention applied to permutation-invariant tabular features, or a consequence of FT-Transformer's specific tokenisation design. Applying the same diagnostic framework to SAINT — which employs dual row-and-column attention — and TabTransformer — which restricts attention to categorical features — would determine whether architectural modifications shift the attention regime and whether such shifts translate into performance improvements on fraud data.

\textbf{Temporal robustness and concept drift.} The BAF dataset preserves the temporal ordering of bank account applications, enabling the construction of time-ordered train/test splits that simulate deployment drift. A direct extension of this work would evaluate whether periodic threshold recalibration alone is sufficient to maintain performance as the distribution evolves, or whether full retraining is necessary — a distinction with significant operational cost implications for fraud detection systems deployed at scale.

\textbf{Ensemble strategies exploiting complementary feature perspectives.} The SHAP analysis establishes that LGBM and FT-Transformer rely on partially overlapping but distinct feature hierarchies (Jaccard\,=\,0.54). This complementarity suggests that a stacking ensemble combining both models' probability scores could potentially exceed either individual model's PR-AUC ceiling — particularly on BAF, where both models reach the same 0.18 limit independently — by leveraging the continuous-feature emphasis of the transformer alongside the categorical-feature dominance of gradient boosting.

\textbf{Native cost-sensitive training for the FT-Transformer.} The current factorial study shows that Class Weights is the safest imbalance intervention for FT-Transformer, but the implementation applies a uniform reweighting to the cross-entropy loss. Asymmetric loss functions designed specifically for extreme imbalance — such as focal loss or direct F$_\beta$ optimisation — integrated natively into the FT-Transformer training loop could provide a more principled alternative that avoids both the distributional distortion of SMOTE and the bluntness of uniform class weighting.

\textbf{Richer feature representations.} The Foggy Vision diagnosis indicates that the current BAF feature space does not provide sufficient discriminative signal to reward selective attention. Augmenting the feature set with derived behavioural features — velocity-based aggregates, graph-based connectivity signals, or device fingerprinting patterns — would test whether FT-Transformer's performance ceiling rises with inputs that better exploit cross-feature attention, the mechanism for which the architecture was designed but which appears underutilised on the current 30-feature tabular representation.
